% Behavior cloning: the naive baseline, its ceiling, and the stitching promise.

\begin{frame}{Why Not Just Copy the Data?}
\begin{itemize}
    \item Simplest possible idea: treat $\mathcal{D}$ as a supervised dataset and
          \emph{imitate} it
    \item \textbf{Behavior cloning} (BC): train a policy $\pi_\theta$ to
          reproduce the data's actions:
          \[ \min_\theta\; -\,\mathbb{E}_{(s,a)\sim\mathcal{D}}
             \big[\log \pi_\theta(a\mid s)\big] \]
    \item In words: at each state in $\mathcal{D}$, push $\pi_\theta$ to put
          \emph{high probability on the action the data actually took}, i.e.\
          $\max_\theta \mathbb{E}[\log \pi_\theta]$. Papers write this as
          ``$\min$ of $-\log$'' (the standard \emph{negative log-likelihood loss}),
          since optimizers minimize losses
\end{itemize}
\end{frame}

\begin{frame}{Why Not Just Copy the Data?}
\begin{itemize}
    \item Simplest possible idea: treat $\mathcal{D}$ as a supervised dataset and
          \emph{imitate} it
    \item \textbf{Behavior cloning} (BC): train a policy $\pi_\theta$ to
          reproduce the data's actions:
          \[ \min_\theta\; -\,\mathbb{E}_{(s,a)\sim\mathcal{D}}
             \big[\log \pi_\theta(a\mid s)\big] \]
    \item In words: at each state in $\mathcal{D}$, push $\pi_\theta$ to put
          \emph{high probability on the action the data actually took}, i.e.\
          $\max_\theta \mathbb{E}[\log \pi_\theta]$. Papers write this as
          ``$\min$ of $-\log$'' (the standard \emph{negative log-likelihood loss}),
          since optimizers minimize losses
    \item Plain supervised learning (like model learning, but for the policy):
          input $s$, target $a$, with no rewards, no Bellman, no bootstrapping.
          What could go wrong?
\end{itemize}
\end{frame}

% === Ceiling and drift (figure) =========================================
\begin{frame}{Behavior Cloning: The Ceiling and the Drift}
\definecolor{oteal}{RGB}{0,141,160}
\definecolor{ogreen}{RGB}{106,168,79}
\definecolor{ored}{RGB}{158,28,38}
\definecolor{ogray}{RGB}{125,125,125}
\centering
\begin{tikzpicture}[font=\scriptsize, >={Stealth[length=2.4mm]}]
    \fill[oteal!14]
      (0,0.35) .. controls (2,1.35) and (4,0.55) .. (7,1.45) --
      (7,2.45) .. controls (4,1.55) and (2,2.35) .. (0,1.35) -- cycle;
    \node[oteal!60!black] at (5.7,2.5) {data support};
    \draw[ogreen!70!black, very thick]
      (0,0.85) .. controls (2,1.85) and (4,1.05) .. (7,1.95)
      node[right=-2pt, font=\scriptsize]{$\pi_\beta$};
    \draw[ored, very thick, dashed, ->]
      (0,0.85) .. controls (1.8,1.75) and (3.2,0.95) ..
      (4.6,0.45) .. controls (5.6,0.1) .. (6.6,-0.3)
      node[below left=-1pt, font=\scriptsize]{$\pi_{\mathrm{BC}}$};
    \draw[->, ored] (2.6,1.42) -- (2.6,1.12)
      node[midway, left, font=\scriptsize]{$\varepsilon$};
    \node[ogray, font=\scriptsize] at (2.0,-0.25) {never-seen states};
\end{tikzpicture}
\begin{itemize}
    \item $\pi_\beta$ = the (unknown) policy that \emph{collected} the data
          ($\beta$ for ``behavior''); $\pi_{\mathrm{BC}}$ = the clone we train
    \item Ceiling: at best $\pi_{\mathrm{BC}}$ \emph{matches} $\pi_\beta$, copying
          its mistakes too
    \item Drift: each step the clone has a small chance $\varepsilon$ of a wrong
          action $\to$ a slightly unfamiliar state $\to$ a bigger error $\to$ off
          the data entirely
    \item These errors \emph{compound}: total error grows like
          $O(\varepsilon T^2)$ over a horizon of $T$ steps (Ross \& Bagnell, 2010)
\end{itemize}
\end{frame}

% === Stitching (figure) =================================================
\begin{frame}{The Promise: Stitching}
\definecolor{oteal}{RGB}{0,141,160}
\definecolor{ogreen}{RGB}{106,168,79}
\definecolor{oorange}{RGB}{230,145,56}
\definecolor{oblue}{RGB}{16,53,103}
\centering
\begin{tikzpicture}[font=\scriptsize, >={Stealth[length=2.4mm]}]
    \coordinate (A) at (0,0);
    \coordinate (B) at (0.4,2.6);
    \coordinate (M) at (3,1.2);
    \coordinate (G) at (6.4,0.6);
    \draw[ogreen, line width=4.5pt, opacity=0.35]
      (A) .. controls (1.2,1.0) and (2.0,1.3) .. (M)
          .. controls (4.2,0.9) and (5.4,0.2) .. (G);
    \draw[->, oorange, very thick]
      (A) .. controls (1.2,1.0) and (2.0,1.3) .. (M)
          .. controls (3.8,1.9) .. (4.8,2.4);
    \node[oorange!80!black, font=\scriptsize] at (5.3,2.55) {wanders off};
    \draw[->, oteal, very thick]
      (B) .. controls (1.2,2.2) and (2.2,1.7) .. (M)
          .. controls (4.2,0.9) and (5.4,0.2) .. (G);
    \fill[oblue] (M) circle (2.2pt);
    \node[oblue, font=\scriptsize, below=2pt of M] {shared state};
    \node[font=\scriptsize, left=1pt] at (A) {traj.\ A};
    \node[font=\scriptsize, left=1pt] at (B) {traj.\ B};
    \node[ogreen!60!black, font=\scriptsize\bfseries] at (6.4,1.25) {goal};
    \node[ogreen!60!black, font=\Large] at (G) {$\star$};
    \node[ogreen!60!black, font=\scriptsize] at (2.0,0.25) {stitched policy};
\end{tikzpicture}
\begin{itemize}
    \item $\mathcal{D}$ may contain \textbf{no good trajectory}, yet contain
          all the \emph{pieces} of one
    \item Trajectories cross at a shared state: dynamic programming can
          \textbf{stitch} the good halves together
    \item The stitched policy is \emph{better than anything in the data};
          BC can never do this
    \item And we already own an off-policy algorithm: Q-learning learns from
          \emph{any} data\ldots{} right?
\end{itemize}
\end{frame}
